% core/preamble.tex — 核心罗塞塔石碑词典 导言
\documentclass[11pt,a4paper]{article}
\usepackage[UTF8, fontset=windows, heading=true]{ctex}
\usepackage{amsmath,amssymb,amsthm}
\usepackage{bm}
\usepackage{geometry}
\geometry{left=2.4cm,right=2.4cm,top=2.6cm,bottom=2.6cm}
\usepackage{xcolor}
\usepackage{graphicx}
\usepackage{enumitem}
\usepackage{booktabs}
\usepackage{array}
\usepackage{longtable}
\usepackage{caption}
\usepackage{tcolorbox}
\tcbuselibrary{skins,breakable}
\usepackage[colorlinks=true,linkcolor=blue!45!black,urlcolor=blue!45!black]{hyperref}
\linespread{1.18}

% ---- 定理类环境（nLab 风格：定义、命题、定理、例、注、猜想）----
\theoremstyle{definition}
\newtheorem{definition}{定义}[section]
\newtheorem{example}[definition]{例}
\newtheorem{remark}[definition]{注}
\theoremstyle{plain}
\newtheorem{theorem}[definition]{定理}
\newtheorem{proposition}[definition]{命题}
\newtheorem{lemma}[definition]{引理}
\newtheorem{conjecture}[definition]{猜想}
\renewcommand{\proofname}{证明要点}

% ---- 颜色 ----
\definecolor{mcol}{RGB}{28,78,140}
\definecolor{pcol}{RGB}{150,74,10}
\definecolor{gcol}{RGB}{105,105,105}

% ---- 词条环境：#1 中文名 #2 English #3 label ----
\newenvironment{centry}[3]{%
  \section[#1]{#1（#2）}%
  \label{core:#3}%
}{%
  \medskip\hrule height 0.3pt\medskip
}

% ---- 七段式小标题 ----
\newcommand{\oneline}[1]{\par\medskip\noindent{\color{mcol}\bfseries 一句话}\quad #1\par}
\newcommand{\physside}[1]{\par\noindent{\color{pcol}\bfseries 物理侧}\quad #1\par}
\newcommand{\mathside}[1]{\par\noindent{\color{mcol}\bfseries 数学侧}\quad #1\par}
\newcommand{\statusline}[1]{\par\noindent{\color{gcol}\bfseries 状态与出处}\quad #1\par}
\newcommand{\dictionaryrow}[1]{\par\noindent{\color{gcol}\bfseries 词典条目}\quad #1\par}

% ---- 石碑总表 ----
\newenvironment{stone}{\par\medskip\centering\small}{\par\medskip}

% ---- 侧栏「余论」----
\newenvironment{yulun}{\begin{tcolorbox}[colback=gray!6,colframe=gcol!55,boxrule=0.4pt,breakable,left=6pt,right=6pt,top=4pt,bottom=4pt]}{\end{tcolorbox}}
